\section{Monitoraggio del link: PDR e ring buffer}

Per poter decidere quando degradare o recuperare lo stato del platoon, \emph{SafeSwitch} ha bisogno di una misura della qualita` del link che sia confrontabile tra tecnologie molto diverse tra loro. Nel paper questa misura e` il \emph{Packet Delivery Ratio} (PDR), cioe` la frazione di pacchetti correttamente ricevuti in una finestra temporale recente \cite{segata2023multi-technology}. L'attenzione è posta sul capire se il flusso informativo disponibile al controllore sia ancora sufficientemente affidabile.

La stima del PDR sfrutta il fatto che ogni beacon venga trasmesso in parallelo sulle interfacce attive. Quando un veicolo riceve piu` copie dello stesso messaggio su tecnologie differenti, puo` confrontarne i numeri di sequenza e ricostruire quali pacchetti siano arrivati e quali invece siano andati persi su ciascun link. In questo modo perdite, interruzioni temporanee o degradi di qualita` vengono osservati non come eventi isolati, ma come andamento nel tempo della consegna dei messaggi.

Per mantenere questa informazione viene usato un \emph{ring buffer}, ovvero una struttura circolare che conserva l'esito degli ultimi messaggi osservati per ciascuna interfaccia. L'aspetto importante, a livello concettuale, non e` la struttura dati in se', ma il fatto che il sistema ragioni su una finestra scorrevole di ricezioni recenti, e non su un singolo evento. Questo permette di filtrare fluttuazioni brevi e di valutare il comportamento del link in modo piu` stabile.

Il monitoraggio viene applicato alle due sorgenti informative piu` rilevanti per il follower: il leader e il predecessore immediato. Le statistiche calcolate sulla finestra recente alimentano poi la logica di \emph{failure} e \emph{recovery}. Un'interfaccia viene considerata degradata quando il suo PDR scende sotto una soglia operativa; al contrario, il recupero richiede che la qualita` torni sopra la soglia e vi rimanga abbastanza a lungo da evitare oscillazioni dovute a recuperi solo transitori.

Un vantaggio di questa impostazione e` che fenomeni anche molto diversi tra loro vengono ricondotti a una stessa astrazione operativa. Contesa del canale 802.11p, perdita della linea di vista per la VLC o transitori di handover sulla rete cellulare hanno cause fisiche differenti, ma per \emph{SafeSwitch} assumono tutti la forma di una riduzione dell'affidabilita` del flusso di messaggi che raggiunge il controllore.
